\chapter{Error-Correcting Codes — MRRW Bound}
%%% generated by the blueprint pipeline from TCSlib.ErrorCorrectingCodes.MRRW
%%%%%%%%%%%%%%%%%%%%%%%%%%%%%%%%%%%%%%%%%%%%%%%%%%%%%%%%%%%%%%%%%%%%%%%%%%%%%%%
\section{Overview}

This chapter formalizes the binary first McEliece--Rodemich--Rumsey--Welch (MRRW)
bound on the asymptotic rate of binary codes: for every $\delta \in [0,1/2]$,
\[
  R(\delta) \;\le\; H\!\left(\tfrac12 - \sqrt{\delta(1-\delta)}\right),
\]
where $R$ is the asymptotic binary code rate and $H$ is the binary entropy function.
The development supplies the maximum code size $A(n,d)$, the binary Krawtchouk
polynomials (both as an explicit alternating sum and as genuine polynomials over
$\mathbb{R}$ defined by the three-term recurrence), Delsarte's linear-programming
bound, the truncated Christoffel--Darboux kernel $\Phi_{n,t}$ together with its
feasibility properties, and the asymptotic ingredients needed to assemble the
main theorem.

%%%%%%%%%%%%%%%%%%%%%%%%%%%%%%%%%%%%%%%%%%%%%%%%%%%%%%%%%%%%%%%%%%%%%%%%%%%%%%%
\section{Declarations}

\subsection{Statement and normalization}

\begin{definition}[Maximum size of a binary code with given distance]
\lean{codeA}
\label{codeA}
\leanok
\statementsource{philips}{
  pdf-24e94ad30862-p025-b003,
  pdf-24e94ad30862-p025-b004
}
\statementsource{IEEE-TIT.23.157}{pdf-abed77da6b7a-p001-b004}
For naturals $n$ and $d$, $A(n,d)$ is the supremum of the cardinalities $k$ for which
there is a finite set $C$ of binary codewords $\mathrm{Fin}\,n \to \mathrm{Bool}$ with
$|C| = k$ and such that any two distinct $x, y \in C$ differ in at least $d$
coordinates.
\end{definition}

\begin{definition}[Binary entropy in base $2$]
\lean{binaryEntropy}
\label{binaryEntropy}
\leanok
\statementsource{IEEE-TIT.23.157}{pdf-abed77da6b7a-p001-b005}
The binary entropy function
\[
  H(x) \;=\; -x\log_2 x - (1-x)\log_2(1-x),
\]
using the base-$2$ logarithm, with the convention $0\log_2 0 = 0$ inherited from
$\log_2 0 = 0$.
\end{definition}

\begin{definition}[Asymptotic binary code rate]
\lean{rate}
\label{rate}
\leanok
\statementsource{IEEE-TIT.23.157}{pdf-abed77da6b7a-p001-b004}
\uses{codeA}
For $\delta \in [0,1]$, the asymptotic rate is
\[
  R(\delta) \;=\; \limsup_{n\to\infty} \frac{1}{n}\log_2 A(n, \lfloor \delta n \rfloor).
\]
\end{definition}

\subsection{Krawtchouk polynomials}

\begin{definition}[Binary Krawtchouk polynomial at integer arguments]
\lean{krawtchouk}
\label{krawtchouk}
\leanok
\statementsource{philips}{pdf-24e94ad30862-p024-b003}
\statementsource{arXiv.1812.04992}{
  pdf-f300deb505d8-p002-b005,
  pdf-f300deb505d8-p002-b006
}
\statementsource{IEEE-TIT.23.157}{pdf-abed77da6b7a-p007-b009}
For naturals $n$, $j$, $x$,
\[
  K_j^{(n)}(x) \;=\; \sum_{i=0}^{j} (-1)^i \binom{x}{i}\binom{n-x}{j-i},
\]
where the subtraction $n-x$ is natural-number subtraction; the definition gives the
intended values for $x \le n$.
\end{definition}

\begin{definition}[Krawtchouk polynomials over $\mathbb{R}$]
\lean{krawtchoukPoly}
\label{krawtchoukPoly}
\leanok
\statementsource{philips}{
  pdf-24e94ad30862-p024-b003,
  pdf-24e94ad30862-p025-b002
}
\statementsource{arXiv.1812.04992}{
  pdf-f300deb505d8-p002-b005,
  pdf-f300deb505d8-p002-b006
}
\statementsource{IEEE-TIT.23.157}{
  pdf-abed77da6b7a-p007-b010,
  pdf-abed77da6b7a-p007-b013
}
The family of real polynomials defined by the three-term recurrence
$K_0 = 1$, $K_1 = n - 2X$, and
\[
  (j+2)\,K_{j+2} \;=\; (n - 2X)\,K_{j+1} - (n-j)\,K_j,
\]
extending the integer-argument definition to real arguments.
\end{definition}

\begin{theorem}[Value of the zeroth Krawtchouk polynomial]
\lean{krawtchouk_zero}
\proofsource{IEEE-TIT.23.157}{pdf-abed77da6b7a-p007-b010}
\label{krawtchouk_zero}
\leanok
\statementsource{IEEE-TIT.23.157}{pdf-abed77da6b7a-p007-b010}

\proofstep
  {krawtchouk}
  {direct}
  {IEEE-TIT.23.157}
  {pdf-abed77da6b7a-p007-b009,
    pdf-abed77da6b7a-p007-b008}
\uses{krawtchouk}
\difficulty{1}
For all $n$ and $x$, $K_0^{(n)}(x) = 1$.
\end{theorem}

\begin{theorem}[Krawtchouk polynomial at $x = 0$]
\lean{krawtchouk_eval_zero}
\statementsource{arXiv.1812.04992}{pdf-f300deb505d8-p003-b007}
\label{krawtchouk_eval_zero}
\leanok
\statementsource{philips}{
  pdf-24e94ad30862-p024-b003,
  pdf-24e94ad30862-p024-b004
}
\statementsource{IEEE-TIT.23.157}{pdf-abed77da6b7a-p007-b010}
\uses{krawtchouk}
\difficulty{2}
For $j \le n$, $K_j^{(n)}(0) = \binom{n}{j}$.
\end{theorem}

\begin{theorem}[First Krawtchouk polynomial]
\lean{krawtchouk_one}
\proofsource{IEEE-TIT.23.157}{pdf-abed77da6b7a-p007-b010}
\label{krawtchouk_one}
\leanok
\statementsource{arXiv.1812.04992}{pdf-f300deb505d8-p002-b006}
\statementsource{IEEE-TIT.23.157}{pdf-abed77da6b7a-p007-b010}

\proofstep
  {krawtchouk}
  {direct}
  {IEEE-TIT.23.157}
  {pdf-abed77da6b7a-p007-b009}
\uses{krawtchouk}
\difficulty{3}
For $x \le n$, $K_1^{(n)}(x) = n - 2x$.
\end{theorem}

\begin{theorem}[Krawtchouk generating function]
\lean{krawtchouk_generating_function}
\label{krawtchouk_generating_function}
\leanok
\statementsource{arXiv.1812.04992}{
  pdf-f300deb505d8-p002-b005,
  pdf-f300deb505d8-p002-b006
}
\statementsource{IEEE-TIT.23.157}{pdf-abed77da6b7a-p007-b008}
\uses{krawtchouk}
\difficulty{6}
For $x \le n$ and every real $z$,
\[
  \sum_{j=0}^{n} K_j^{(n)}(x)\, z^{j} \;=\; (1-z)^{x}(1+z)^{\,n-x}.
\]
\end{theorem}

\begin{theorem}[Orthogonality of Krawtchouk polynomials]
\lean{krawtchouk_orthogonality}
\label{krawtchouk_orthogonality}
\leanok
\statementsource{philips}{pdf-24e94ad30862-p025-b002}
\statementsource{arXiv.1812.04992}{pdf-f300deb505d8-p002-b005}
\statementsource{IEEE-TIT.23.157}{pdf-abed77da6b7a-p007-b012}
\uses{krawtchouk, krawtchouk_generating_function}
\difficulty{7}
For $r \le n$ and $s \le n$,
\[
  \sum_{x=0}^{n} \binom{n}{x} K_r^{(n)}(x)\,K_s^{(n)}(x)
  \;=\;
  \begin{cases}
    2^{n}\binom{n}{r}, & r = s,\\
    0, & r \ne s.
  \end{cases}
\]
\end{theorem}

\begin{theorem}[Three-term recurrence for Krawtchouk polynomials]
\lean{krawtchouk_recurrence}
\label{krawtchouk_recurrence}
\leanok
\statementsource{philips}{pdf-24e94ad30862-p025-b002}
\statementsource{arXiv.1812.04992}{pdf-f300deb505d8-p002-b005}
\statementsource{IEEE-TIT.23.157}{pdf-abed77da6b7a-p007-b013}
\uses{krawtchouk, krawtchouk_generating_function}
\difficulty{7}
For $1 \le j$, $j + 1 \le n$ and $x \le n$,
\[
  (j+1)\,K_{j+1}^{(n)}(x) \;=\; (n - 2x)\,K_j^{(n)}(x) - (n - j + 1)\,K_{j-1}^{(n)}(x).
\]
\end{theorem}

\begin{theorem}[Agreement of the sum formula and the polynomial form]
\lean{krawtchoukPoly_eval_nat}
\label{krawtchoukPoly_eval_nat}
\leanok
\uses{krawtchouk, krawtchoukPoly, krawtchouk_zero, krawtchouk_one, krawtchouk_recurrence}
\difficulty{4}
For $j \le n$ and $x \le n$, evaluating the polynomial $K_j^{(n)}$ at the real number
$x$ gives the value $K_j^{(n)}(x)$ of the alternating-sum definition.
\end{theorem}

\subsection{Delsarte's linear programming bound}

\begin{theorem}[Delsarte LP inequality]
\lean{delsarte_lp_bound}
\proofsource{IEEE-TIT.23.157}{
  pdf-abed77da6b7a-p005-b002,
  pdf-abed77da6b7a-p005-b003
}
\label{delsarte_lp_bound}
\statementsource{philips}{
  pdf-24e94ad30862-p018-b006,
  pdf-24e94ad30862-p019-b002,
  pdf-24e94ad30862-p020-b006,
  pdf-24e94ad30862-p024-b006,
  pdf-24e94ad30862-p025-b003
}
\statementsource{IEEE-TIT.23.157}{
  pdf-abed77da6b7a-p003-b002,
  pdf-abed77da6b7a-p005-b002,
  pdf-abed77da6b7a-p005-b003
}

\proofstep
  {codeA}
  {context}
  {IEEE-TIT.23.157}
  {pdf-abed77da6b7a-p001-b004}
\proofstep
  {krawtchouk}
  {context}
  {IEEE-TIT.23.157}
  {pdf-abed77da6b7a-p007-b009,
    pdf-abed77da6b7a-p007-b008}
\uses{codeA, krawtchouk}
Let $F(x) = \sum_{j=0}^{n} F_j K_j^{(n)}(x)$ for real coefficients $F_j$. If
$F_0 > 0$, if $F_j \ge 0$ for all $1 \le j \le n$, and if $F(x) \le 0$ for every
integer $x$ with $d \le x \le n$, then
\[
  A(n,d) \;\le\; \frac{F(0)}{F_0}.
\]
\end{theorem}

\subsection{Christoffel--Darboux kernel and feasibility}

\begin{definition}[Truncated Christoffel--Darboux kernel]
\lean{cdKernel}
\label{cdKernel}
\leanok
\statementsource{IEEE-TIT.23.157}{pdf-abed77da6b7a-p005-b004}
\uses{krawtchoukPoly}
For $t \le n$ and real $a, x$, the truncated reproducing kernel is
\[
  \Phi_{n,t}(a,x) \;=\; \sum_{j=0}^{t} \frac{K_j^{(n)}(a)\,K_j^{(n)}(x)}{\binom{n}{j}}.
\]
\end{definition}

\begin{theorem}[Christoffel--Darboux identity for Krawtchouk polynomials]
\lean{cd_identity}
\label{cd_identity}
\leanok
\statementsource{IEEE-TIT.23.157}{
  pdf-abed77da6b7a-p008-b002,
  pdf-abed77da6b7a-p008-b003,
  pdf-abed77da6b7a-p005-b004
}
\uses{cdKernel, krawtchoukPoly}
\difficulty{6}
For $t + 1 \le n$ and real $a \ne x$ there is a scalar $c > 0$ with
\[
  \Phi_{n,t}(a,x) \;=\; c\,\frac{K_t^{(n)}(a)K_{t+1}^{(n)}(x) - K_{t+1}^{(n)}(a)K_t^{(n)}(x)}{a - x}.
\]
\end{theorem}

\begin{theorem}[Positivity of Krawtchouk values below all zeros]
\lean{krawtchouk_positive_below_smallest_zero}
\proofsource{IEEE-TIT.23.157}{
  pdf-abed77da6b7a-p005-b005,
  pdf-abed77da6b7a-p008-b004
}
\label{krawtchouk_positive_below_smallest_zero}
\statementsource{arXiv.1812.04992}{
  pdf-f300deb505d8-p003-b005,
  pdf-f300deb505d8-p003-b007
}
\statementsource{IEEE-TIT.23.157}{
  pdf-abed77da6b7a-p005-b005,
  pdf-abed77da6b7a-p008-b004
}

\proofstep
  {krawtchoukPoly}
  {direct}
  {IEEE-TIT.23.157}
  {pdf-abed77da6b7a-p007-b013,
    pdf-abed77da6b7a-p007-b010,
    pdf-abed77da6b7a-p007-b008}
\uses{krawtchoukPoly}
Fix $t + 1 \le n$ and a real $a$ that is strictly smaller than every real zero
$\xi$ of $K_t^{(n)}$. Then $K_j^{(n)}(a) > 0$ for every $j \le t$.
\end{theorem}

\begin{theorem}[Nonpositivity of the kernel beyond a threshold]
\lean{cdKernel_nonpos_beyond_threshold}
\label{cdKernel_nonpos_beyond_threshold}
\statementsource{IEEE-TIT.23.157}{
  pdf-abed77da6b7a-p005-b004,
  pdf-abed77da6b7a-p005-b005
}
\uses{cdKernel, krawtchoukPoly}
Under the same hypotheses ($t + 1 \le n$ and $a$ strictly below every real zero of
$K_t^{(n)}$), there exists a threshold $\eta$ such that $\Phi_{n,t}(a,x) \le 0$ for
every integer $x$ with $\eta \le x \le n$.
\end{theorem}

\begin{theorem}[Finite-$n$ MRRW bound]
\lean{finite_n_mrrw_bound}
\label{finite_n_mrrw_bound}
\leanok
\statementsource{IEEE-TIT.23.157}{
  pdf-abed77da6b7a-p005-b004,
  pdf-abed77da6b7a-p005-b005
}
\uses{cdKernel, codeA, delsarte_lp_bound, krawtchoukPoly, krawtchoukPoly_eval_nat, krawtchouk_positive_below_smallest_zero}
\difficulty{6}
Let $t + 1 \le n$, let $a$ be strictly below every real zero of $K_t^{(n)}$, and
suppose $\Phi_{n,t}(a,x) \le 0$ for every integer $x$ with $d \le x \le n$. Then
\[
  A(n,d) \;\le\; \Phi_{n,t}(a,0).
\]
\end{theorem}

\subsection{Asymptotic choice of parameters}

\begin{theorem}[Asymptotics of the smallest Krawtchouk zero]
\lean{smallest_zero_asymptotic}
\proofsource{IEEE-TIT.23.157}{
  pdf-abed77da6b7a-p008-b007,
  pdf-abed77da6b7a-p008-b008
}
\label{smallest_zero_asymptotic}
\statementsource{arXiv.1812.04992}{
  pdf-f300deb505d8-p004-b007,
  pdf-f300deb505d8-p005-b001
}
\statementsource{IEEE-TIT.23.157}{
  pdf-abed77da6b7a-p008-b007,
  pdf-abed77da6b7a-p008-b008,
  pdf-abed77da6b7a-p006-b002
}

\proofstep
  {krawtchoukPoly}
  {direct}
  {IEEE-TIT.23.157}
  {pdf-abed77da6b7a-p007-b013,
    pdf-abed77da6b7a-p007-b010,
    pdf-abed77da6b7a-p007-b008,
    pdf-abed77da6b7a-p007-b007}
\proofstep
  {cdKernel}
  {direct}
  {IEEE-TIT.23.157}
  {pdf-abed77da6b7a-p005-b004,
    pdf-abed77da6b7a-p008-b003,
    pdf-abed77da6b7a-p008-b002}
\uses{cdKernel, krawtchoukPoly}
Let $0 < \tau < 1/2$ and let $t_n$ satisfy $t_n / n \to \tau$. Then there is a
sequence $a_n$ such that, eventually in $n$, $a_n$ lies strictly below every real
zero of $K_{t_n}^{(n)}$, with $a_n/n \to \tfrac12 - \sqrt{\tau(1-\tau)}$, and such
that eventually $\Phi_{n,t_n}(a_n,x) \le 0$ for all integers $x$ with
$\lfloor (\tfrac12 - \sqrt{\tau(1-\tau)})n \rfloor \le x \le n$.
\end{theorem}

\begin{theorem}[Entropy growth rate of the objective]
\lean{entropy_growth_of_objective}
\proofsource{IEEE-TIT.23.157}{
  pdf-abed77da6b7a-p006-b002,
  pdf-abed77da6b7a-p006-b001
}
\label{entropy_growth_of_objective}
\statementsource{IEEE-TIT.23.157}{
  pdf-abed77da6b7a-p006-b001,
  pdf-abed77da6b7a-p006-b002
}

\proofstep
  {binaryEntropy}
  {direct}
  {IEEE-TIT.23.157}
  {pdf-abed77da6b7a-p001-b005}
\proofstep
  {krawtchoukPoly}
  {direct}
  {IEEE-TIT.23.157}
  {pdf-abed77da6b7a-p007-b013,
    pdf-abed77da6b7a-p007-b008,
    pdf-abed77da6b7a-p007-b010}
\proofstep
  {cdKernel}
  {direct}
  {IEEE-TIT.23.157}
  {pdf-abed77da6b7a-p008-b002,
    pdf-abed77da6b7a-p008-b003,
    pdf-abed77da6b7a-p005-b004}
\uses{binaryEntropy, cdKernel, krawtchoukPoly}
Let $0 < \tau < 1/2$, let $t_n/n \to \tau$, and let $a_n$ eventually lie strictly
below every real zero of $K_{t_n}^{(n)}$. Then
\[
  \limsup_{n\to\infty} \frac{1}{n}\log_2 \Phi_{n,t_n}(a_n,0) \;\le\; H(\tau).
\]
\end{theorem}

\begin{theorem}[Entropy asymptotics for binomial tails]
\lean{binomial_tail_entropy_asymptotic}
\label{binomial_tail_entropy_asymptotic}
\uses{binaryEntropy}
For $0 \le \tau \le 1/2$,
\[
  \frac{1}{n}\log_2\!\left(\sum_{j=0}^{\lfloor \tau n\rfloor} \binom{n}{j}\right)
  \;\longrightarrow\; H(\tau)
  \qquad (n \to \infty).
\]
\end{theorem}

\subsection{Conclusion of the MRRW bound}

\begin{theorem}[Involution between $\delta$ and $\tau$]
\lean{delta_tau_involution}
\label{delta_tau_involution}
\leanok
\difficulty{3}
For $0 \le \delta \le 1/2$, setting $\tau = \tfrac12 - \sqrt{\delta(1-\delta)}$ gives
back $\delta = \tfrac12 - \sqrt{\tau(1-\tau)}$; that is, the map is an involution on
$[0,1/2]$.
\end{theorem}

\begin{theorem}[Binary entropy at $0$]
\lean{binaryEntropy_zero}
\label{binaryEntropy_zero}
\leanok
\uses{binaryEntropy}
\difficulty{1}
$H(0) = 0$.
\end{theorem}

\begin{theorem}[Binary entropy at $1/2$]
\lean{binaryEntropy_half}
\label{binaryEntropy_half}
\leanok
\uses{binaryEntropy}
\difficulty{2}
$H(1/2) = 1$.
\end{theorem}

\begin{theorem}[Binary first MRRW bound]
\lean{mrrw_bound}
\proofsource{IEEE-TIT.23.157}{
  pdf-abed77da6b7a-p003-b003,
  pdf-abed77da6b7a-p006-b002
}
\label{mrrw_bound}
\statementsource{IEEE-TIT.23.157}{
  pdf-abed77da6b7a-p001-b006,
  pdf-abed77da6b7a-p003-b003,
  pdf-abed77da6b7a-p006-b002
}

\proofstep
  {codeA}
  {direct}
  {IEEE-TIT.23.157}
  {pdf-abed77da6b7a-p001-b004}
\proofstep
  {binaryEntropy}
  {direct}
  {IEEE-TIT.23.157}
  {pdf-abed77da6b7a-p001-b005}
\proofstep
  {rate}
  {direct}
  {IEEE-TIT.23.157}
  {pdf-abed77da6b7a-p001-b004}
\uses{binaryEntropy, rate}
For every $\delta$ with $0 \le \delta \le 1/2$,
\[
  R(\delta) \;\le\; H\!\left(\tfrac12 - \sqrt{\delta(1-\delta)}\right).
\]
\end{theorem}
